\chapter{Driving the junction}\label{ch:driving}

\section{Which parameters can be driven}\label{sec:drives}

In \pkg{} a drive is a periodic modulation of one or several model parameters,
\begin{equation}\label{eq:drive}
  p(t)=p_0+\sum_n A_n\cos(n\Omega t+\theta_n),
\end{equation}
and the driven Hamiltonian is $H(t)=H\bigl(p(t)\bigr)$, built by the same model
constructor as the static one. For a monochromatic drive we write $A\equiv A_1$ and
$\theta\equiv\theta_1$. To distinguish the driven parameters, $A_\phi$ denotes the amplitude
of a phase drive (in radians) and $A_\eps$ that of a gate drive (an energy). The physically
important cases are the following.

\begin{description}[leftmargin=0pt]
\item[Gate drive.] An ac voltage on a gate electrode modulates the dot level,
$\eps(t)$. Because $\Hd$ is linear in $\eps$, $H(t)=H(\eps_0)+(n_d-1)\sum_nA_n\cos(n\Omega
t+\theta_n)$. A monochromatic drive has only the harmonics
$H_{\pm1}=\tfrac{A_\eps}2\ee^{\pm\ii\theta}(n_d-1)$.
This is the setting of photon-assisted tunnelling through quantum
dots~\cite{kouwenhoven1994photon,kouwenhoven1994observation,oosterkamp1998microwave,platero2004photon}.
\item[Phase drive.] Microwaves applied to a superconducting loop or across the junction
modulate the phase difference, $\phi(t)$. Since the phase enters the tunnelling through
the factors $\ee^{\pm\ii\phi/4}$, $H(t)$ contains all harmonics, with Bessel-function
amplitudes (\cref{sec:jacobi-anger}). This is how Andreev levels are
excited in microwave spectroscopy~\cite{bretheau2013exciting,janvier2015coherent,hays2018direct,kos2013frequencydependent}.
\item[Coupling drive.] A gate on a tunnel barrier modulates $\Gamma_l$, which enters the
Hamiltonian through $t_l\propto\sqrt{\Gamma_l}$~\cite{moghaddam2012driven}.
\item[Field drive.] An oscillating magnetic field along $z$ modulates the Zeeman energy. A
transverse field $b(t)S^x_d$ flips the spin and breaks the conservation of $S_z$; this is
electron spin resonance (ESR) of the impurity spin
(\cref{ch:ex-esr})~\cite{koppens2006driven,baumann2015electron,yang2019coherent}.
\item[Combined drives.] Several parameters may be modulated simultaneously with relative
phase offsets, for example by a gate electrode that also couples capacitively to the junction phase.
\end{description}

\begin{incode}
\begin{lstlisting}[frame=none,backgroundcolor={},aboveskip=0pt,belowskip=0pt]
from qdjj_solver import Harmonic, ParameterDrive, PeriodicHamiltonian, reference_model
from qdjj_solver.qp_solver import solve_floquet
drives = [ParameterDrive.cosine("phi", .3),                      # phi0 + 0.3 cos(Omega t)
          ParameterDrive("detuning", (Harmonic(1, .02, np.pi/2),  # two harmonics
                                      Harmonic(2, .01)))]
periodic = PeriodicHamiltonian.from_parameters(reference_model, model, drives, omega=.5)
\end{lstlisting}
The model constructor, here \code{reference_model}, is evaluated at equally spaced times
over one period (\code{samples}, default 64), and every operator coefficient is Fourier
transformed (\cref{sec:sampling}). \code{from_callable(f, omega)} accepts any function
\code{f(t)} that returns a \code{Hamiltonian}.
\par\smallskip\usermanual{docs/floquet.md}.
\end{incode}

\section{Phase drive, voltage and the Josephson relation}\label{sec:phase-voltage}

The superconducting phase is included in the tunnelling amplitudes,
\cref{eq:tunnelling}. With this choice a time-dependent phase is equivalent to a voltage.
The periodic transformation $W(t)=\exp[-\ii\sum_l\phi_l(t)N_l/2]$, with $N_l$ the electron
number of lead $l$, removes the phases from the tunnelling and moves them onto the pair
potentials. By \cref{eq:gauge} it also adds the term $\sum_l(\hbar\dot\phi_l/2)N_l$, which
raises the electrochemical potential $\mu_l$ of lead $l$ by $\hbar\dot\phi_l/2$. We define the
bias voltage $V=V_L-V_R$ as the electric potential of the left lead relative to the right
one. Since electrons carry the charge $-e$, $eV=\mu_R-\mu_L$, and the phase difference
$\phi=\phi_R-\phi_L$ obeys the Josephson relation
\begin{equation}\label{eq:josephson}
  \dot\phi=\frac{2e}{\hbar}V .
\end{equation}
With the current $I$ of \cref{eq:current}, positive from the left to the right lead, the
power supplied by the voltage source is $IV=\dot\phi\,\partial E/\partial\phi$. For a single
lead, the phase can instead be moved onto the dot with $W(t)=\ee^{\ii\phi(t)n_d/2}$, which
gives
\begin{equation}\label{eq:phase-to-level}
  H'(t)=H(\phi=0)-\frac{\dot\phi(t)}2\,n_d ,
\end{equation}
i.e.\ a level drive $\eps(t)=\eps_0-\dot\phi/2$ with identical quasienergies. The code is
tested against this identity (\cref{sec:validation-driven}). The same gauge freedom
underlies the description of voltage-biased junctions and multiple Andreev reflection as
Floquet scattering problems~\cite{averin1995ac,bratus1995theory,cuevas1996hamiltonian,cuevas2002subharmonic}.

\section{Harmonics of the phase drive}\label{sec:jacobi-anger}

For the reference junction with $\phi(t)=\phi_0+A\cos(\Omega t+\theta)$, where $A\equiv A_\phi$
in this section, each tunnelling factor $\ee^{\pm\ii\phi(t)/4}$ contains
$\ee^{\pm\ii(A/4)\cos(\Omega t+\theta)}$. The Jacobi--Anger expansion
$\ee^{\ii z\cos\vartheta}=\sum_n\ii^nJ_n(z)\ee^{\ii n\vartheta}$~\cite{olver2026nist}
gives the harmonics in closed form (\cref{app:jacobi}):
\begin{align}
  H_0&=H(\phi_0)\big\vert_{\Gamma\to\Gamma J_0(A/4)^2},\label{eq:phase-h0}\\
  H_k&=4(-1)^{(k-1)/2}J_k\!\left(\tfrac A4\right)\ee^{\ii k\theta}\,\frac{\partial H}{\partial\phi}\bigg\vert_{\phi_0}
  &&(k\ \mathrm{odd}),\label{eq:phase-hk}\\
  H_k&=(-1)^{k/2}J_k\!\left(\tfrac A4\right)\ee^{\ii k\theta}\,H_T(\phi_0)&&(k\ne0\ \mathrm{even}),
\end{align}
where $H_T$ is the tunnelling Hamiltonian. Three consequences are useful. First, for a weak drive
$H_{\pm1}\approx\tfrac A2\ee^{\pm\ii\theta}\partial H/\partial\phi$: the phase drive
couples through the current operator, and the one-photon Rabi splitting between two states
is $A\abs{\matel b{\partial H/\partial\phi}a}$. Second, the time-averaged hybridization is
reduced by the factor $J_0(A/4)^2$, which vanishes at $A\approx9.62$ (\cref{ch:ex-highfreq}).
Third, the even harmonics are proportional to $H_T$ and the odd ones to
$\partial H/\partial\phi$, so two-photon processes probe a different operator from
one-photon processes.

\section{From laboratory parameters to drive amplitudes}\label{sec:lab-units}

The code works with energies in units of $\Delta$ and dimensionless phases. Typical
conversions are:
\begin{itemize}
\item \emph{Frequency.} $\Omega/\Delta=hf/\Delta$; for an aluminium gap $\Delta\approx
180\,\mu$eV, $\Omega=1$ corresponds to $f\approx44$\,GHz, and a typical microwave frequency
of $10$\,GHz to $\Omega\approx0.23$.
\item \emph{Voltage drive.} An ac voltage $V_{\mathrm{ac}}\cos\Omega t$ across the junction
gives $\dot\phi=(2e/\hbar)V_{\mathrm{ac}}\cos\Omega t$, hence the phase amplitude
$A_\phi=2eV_{\mathrm{ac}}/(\hbar\Omega)$, the familiar Tien--Gordon modulation index of Cooper
pairs.
\item \emph{Flux drive.} In a loop threaded by a flux $\Phi$, $\phi=2\pi\Phi/\Phi_0$ with
$\Phi_0=h/2e$, so a flux modulation $\delta\Phi$ gives $A_\phi=2\pi\,\delta\Phi/\Phi_0$.
\item \emph{Gate drive.} An ac gate voltage of amplitude $V_g$ shifts the level by
$\alpha eV_g$, with the lever arm $\alpha$; hence $A_\eps=\alpha eV_g/\Delta$.
\end{itemize}
Microwave spectroscopy usually operates in the weak-drive regime, $A_\phi\ll1$ and
$A_\eps\ll\Omega$. Floquet engineering and the strong-driving phenomena of
\cref{ch:ex-tls,ch:ex-highfreq} require modulation indices of order one.

\section{Fixed canonical coordinates}\label{sec:fixed-coordinates}

The QP basis is built from the Bogoliubov quasiparticles of the isolated reservoirs, which
depend on $\Delta_l$ and on the bath levels $\xi_{li}$ and remain fixed as $\eps$, $U$, $B$,
$\Gamma_l$ or $\phi_l$ vary. All drives of the latter parameters therefore act on a fixed Fock
space, and the Sambe construction applies directly. If the coordinates themselves varied
with time, e.g.\ for a modulated gap $\Delta(t)$, the transformed Hamiltonian would acquire
the additional term $\ii(\partial_tW)W\dg$ of \cref{eq:gauge}. The implemented methods
use fixed coordinates: \pkg{} checks that those of every sampled Hamiltonian
coincide with those of the undriven model, and rejects drives that change them. The eta
representation and paired-mode compression depend on $\eps$ and $\phi$, so driven models
are built with \code{symmetry=False, compress=False}.

\section{Selection rules}\label{sec:selection}

A drive conserves fermion parity, so it couples only states of the same parity.
Parity-changing transitions require the exchange of a quasiparticle with an external
system. Drives that conserve $S_z$ (gate, phase, coupling, longitudinal field) act within
sectors of fixed $S_z$, and the calculation is done sector by sector. A transverse field
requires \code{Sector(parity, None)}.

At the particle--hole-symmetric point of the symmetric junction the eta parity provides a
further, less obvious selection rule. The phase drive and the coupling drive are eta-even and
couple only states of equal eta parity; the gate drive, proportional to $n_d-1$, is eta-odd
and couples only states of opposite eta parity. The lowest two even states $S$ and $S^*$
have opposite eta parity, so at $\eps=0$ the pair transition $S\to S^*$ is invisible to the
phase drive and maximally visible to the gate drive. Away from $\eps=0$ both drives couple
the pair. \Cref{ch:ex-gate} demonstrates these complementary rules. For a non-interacting
resonant level the same rule reads as follows: the symmetric junction at resonance is
reflectionless, and the current matrix element of the pair transition, proportional to
$\sqrt{1-\tau}$ with $\tau$ the transmission, vanishes~\cite{zazunov2003andreev,kos2013frequencydependent}.

\section{The driven proximitized dot}\label{sec:driven-largegap}

The large-gap limit \eqref{eq:largegap} gives a transparent model of the driven even sector.
We write the pseudospin Hamiltonian as
\begin{equation}\label{eq:pseudospin}
  H=\frac U2+\eps\,\tau_z-\abs{\Gamma_\phi}\,\tau_x
   =\frac U2+R\,(\cos\vartheta\,\tau_z-\sin\vartheta\,\tau_x),\qquad
  R=\sqrt{\eps^2+\abs{\Gamma_\phi}^2},\quad \tan\vartheta=\frac{\abs{\Gamma_\phi}}{\eps}.
\end{equation}
The level splitting is $\omega_0=E_{S^*}-E_S=2R$. A gate drive $A_\eps\cos\Omega t\,\tau_z$ has
the transverse component $A_\eps\sin\vartheta$ in the eigenbasis. A phase drive modulates
$\Gamma_\phi=\Gamma\cos(\phi/2)$ and has the transverse component
$\tfrac{A_\phi}2\Gamma\abs{\sin(\phi_0/2)}\cos\vartheta$. The resonant Rabi splittings are
therefore
\begin{equation}\label{eq:largegap-rabi}
  \delta_{\mathrm{gate}}=A_\eps\abs{\sin\vartheta}=\frac{A_\eps\abs{\Gamma_\phi}}{R},\qquad
  \delta_{\mathrm{phase}}=\frac{A_\phi\Gamma}2\abs{\sin\frac{\phi_0}2}\abs{\cos\vartheta}
  =\frac{A_\phi\Gamma}{2}\abs{\sin\frac{\phi_0}2}\frac{\abs\eps}{R}.
\end{equation}
The phase-drive splitting vanishes at $\eps=0$ and the gate-drive splitting is largest
there: the selection rule of \cref{sec:selection} in its simplest form. For a junction with
finite $\Delta$, these expressions provide qualitative guidance; quantitative results use
the numerically calculated matrix elements.

\begin{caution}
The large-gap model assumes that the drive is slow compared with $\Delta$. A strong phase
drive that is fast compared with the dot dynamics but slow compared with $\Delta$
($\Omega\ll\Delta$) renormalizes $\Gamma_\phi$ by the time average
$\langle\cos(\phi(t)/2)\rangle=J_0(A_\phi/2)\cos(\phi_0/2)$. A drive faster than the
bandwidth instead renormalizes each hopping separately by $J_0(A_\phi/4)$, hence
$\Gamma\to\Gamma J_0(A_\phi/4)^2$. The order of the large-gap and high-frequency limits
therefore determines the effective coupling.
\end{caution}
